% Fig. 3 body: the single-entry thermodynamic cycle of theory/derivation_body.tex Section 7
% (lane p0-w3; the commit drawn from is recorded in figures/tikz/schematics.yaml). Included by
% fig3.tex (default, no text) and variants/fig3_symbols.tex (\symbolstrue: symbols on arrows,
% built only on request for the SI). Every state and arrow is identified by position in
% figures/src/captions/fig3.md.
%
% Layout (P0-W6.3): one row of three states, left to right, with straight horizontal arrows
% between them. Every arrow of the cycle (style proc) is a chain of horizontal and vertical
% segments with right-angle bends, written as absolute coordinates joined by `--` (the tests
% parse this). Only the three force vectors at the cap edge (style force) keep their physical
% directions; both styles have the same line weight.
%
% States and arrows taken from Section 7 (no step added):
%   A  bulk reservoir: hydrated Ca2+ and CO3 2- at the regime activities (bulk chemical potentials)
%   B  template-liquid interface at the regime Ca2+ activity: Ca2+ adsorbed at catecholate and
%      carboxylate sites; its free energy is gamma_sl = gamma_sl0 + Delta_gamma_ads (Gibbs-Langmuir)
%   C  cap on the template (base 13.0..16.0 cm, contact angle 55 deg: partial wetting):
%      template-nucleus contact (gamma_sn, a grand-canonical excess) with Ca2+ bridging catecholate
%      and carboxylate sites to the crystal (the contact excess Gamma_Ca,sn, inside gamma_sn),
%      nucleus-liquid surface (gamma_nl), template-liquid interface outside the contact (gamma_sl);
%      Young at the contact line
%   1  A -> B  Ca2+ from the hydrated bulk to a site, Delta_G_seq; enters only gamma_sl (blue)
%   2  B -> C  the cap replaces template-liquid interface by template-nucleus contact over its base
%   3  A -> C  carbonate route along the top: the formula units of the nucleus come from the
%              reservoir at the bulk chemical potentials (driving force k_B T ln Omega_bulk); none
%              comes from the bound layer, and the route enters the cap from above
\def\b{0.2}        % pictogram bond length (cm)
\def\yt{1.6}       % template-liquid interface
\def\ybase{1.0}    % oxide top / template bottom
\begin{tikzpicture}[sharp corners]
\useasboundingbox (0,0) rectangle (17.1,4.4);

% ---------------------------------------------------------------- state boxes (one row)
\fill[statebox] (0.15,0.3) rectangle (5.05,3.35);
\fill[statebox] (6.1,0.3) rectangle (11.0,3.35);
\fill[statebox] (12.05,0.3) rectangle (16.95,3.35);

% ---------------------------------------------------------------- A: bulk reservoir
\hydratedca{1.75,1.85}{0.26}
\pic[scale=0.24] at (3.9,1.85) {carbonate};

% ---------------------------------------------------------------- B: adsorbed template-liquid interface
\oxidestrip{6.25}{10.85}{\ybase}
\fill[template] (6.25,\ybase) rectangle (10.85,\yt);
\foreach \x in {6.75,8.0,9.25} {\pic[scale=\b] at (\x,1.35) {catecholatedianion};}
\pic[scale=\b] at (10.35,1.56) {acetate};
\draw[bold, draw=okblue] (6.25,\yt) -- (10.85,\yt);
\caion{8.0,1.83}
\caion{10.35,1.83}

% ---------------------------------------------------------------- C: cap on the template
\oxidestrip{12.2}{16.8}{\ybase}
\fill[template] (12.2,\ybase) rectangle (16.8,\yt);
\foreach \x in {12.55,13.8,15.2} {\pic[scale=\b] at (\x,1.35) {catecholatedianion};}
\pic[scale=\b] at (14.5,1.56) {acetate};
\pic[scale=\b] at (16.45,1.56) {acetate};
% nucleus (amorphous CaCO3 cap): base 13.0..16.0, contact angle 55 deg (radius 1.8312, centre y 0.5497)
\fill[nucleus] (16.0,\yt) arc[start angle=35, end angle=145, radius=1.8312] -- cycle;
\pic[scale=0.17, rotate=20] at (13.95,2.03) {carbonate};
\pic[scale=0.17, rotate=-15] at (15.05,2.03) {carbonate};
\filldraw[ion] (14.5,2.12) circle (0.11);
\filldraw[ion] (13.45,1.86) circle (0.11);
\filldraw[ion] (15.55,1.86) circle (0.11);
% interfaces: template-liquid outside the contact, template-nucleus contact, nucleus-liquid surface
\draw[bold, draw=okblue] (12.2,\yt) -- (13.0,\yt);
\draw[bold, draw=okblue] (16.0,\yt) -- (16.8,\yt);
\draw[bold, draw=okvermillion] (13.0,\yt) -- (16.0,\yt);
\draw[bold, draw=okblack] (16.0,\yt) arc[start angle=35, end angle=145, radius=1.8312];
% ion-bridged Ca2+ at the sites under the contact: the contact excess Gamma_Ca,sn, inside gamma_sn
\foreach \x in {13.8,14.5,15.2} {\filldraw[ion] (\x,1.72) circle (0.11);}
% Ca2+ adsorbed at the template-liquid interface outside the contact
\caion{12.55,1.83}
\caion{16.45,1.83}
% Young's relation at the right contact line: the three interfacial energies as line tensions,
% in their physical directions (gamma_nl along the cap tangent, 180 - 55 = 125 deg)
\draw[force, draw=okblue] (16.0,\yt) -- (16.75,\yt);
\draw[force, draw=okvermillion] (16.0,\yt) -- (15.3,\yt);
\draw[force, draw=okblack] (16.0,\yt) -- ++(125:0.8);

% ---------------------------------------------------------------- arrows of the cycle (orthogonal)
\draw[proc, draw=okblue] (5.15,2.75) -- (6.0,2.75);                                % 1: Delta_G_seq
\draw[proc] (11.1,2.75) -- (11.95,2.75);                                           % 2: cap formation
\draw[proc] (3.9,2.5) -- (3.9,3.95) -- (14.5,3.95) -- (14.5,2.55);                 % 3: carbonate route to the nucleus

\ifsymbols
  \begingroup\fontsize{7}{8.4}\selectfont
  \node[okblue, anchor=south] at (5.575,2.8) {\sym{Δ\textit{G}}{seq}};
  \node[okblue, anchor=north] at (16.45,\yt-0.08) {\sym{\textit{γ}}{sl}};
  \node[okvermillion, anchor=north] at (15.55,\yt-0.08) {\sym{\textit{γ}}{sn}};
  \node[okblack, anchor=east] at ($(16.0,\yt)+(125:0.8)$) {\sym{\textit{γ}}{nl}};
  \endgroup
\fi
\end{tikzpicture}
